%!TEX root = ../../springer-mono-format/book.tex
%!TEX root = ../../thesis-trier.tex


\IN{commutative automaton}
An automaton $\semiautomat{A}{Q}{\Sigma}{\delta}$ is \emph{commutative} 
if $\applytrans{\delta}{q}{ab} = \applytrans{\delta}{q}{ba}$ for all $q \in Q$
and $a,b \in \Sigma$.
Said differently, the mapping $q \mapsto \applytrans{\delta}{q}{x}$ 
is a morphism for $x \in \Sigma$ 
and so an automaton is commutative precisely if its transformation
semigroup is contained in its endomorphism semigroup.
Thus, the letters fulfill
the properties from \cref{prelim:sec:endo}.
For example, by \cref{prelim:eqn:morph-reachability-order},		
we find that \cref{constr_sync:eqn:reach_compatible}
holds, and so commutative automata are $\rho$-automata, 
$q = \delta(q, u)$ for some $q \in Q$
and $u \in \Sigma^*$
implies $\delta(q, v) = \delta(\delta(q, v), u)$
for all $v \in \Sigma^*$,
and, by \cref{prelim:automata:cor:SCC-disjoint-or-permutates},
every strongly connected component is a block.
Every unary automaton is commutative.
See \cref{fig:ex_comm_aut} for two commutative example automata.
Clearly, for commutative $\mathcal A$,
$L(\mathcal A_{q,F})$ is commutative for each $q \in Q$
and $F \subseteq Q$.

First, \cref{prelim:classes:comm:subsec:rees} we characterize the commutativity of Rees quotient automata.
Then in \cref{prelim:classes:comm:subsec:SCC} we investigate strongly connected commutative automata
and give structural results on them and their endomorphisms. \Cref{prelim:classes:comm:subsec:subset} gives
some results on subset mapping and synchronization properties in commutative automata.
Lastly, \cref{prelim:classes:comm:subsec:extremal} characterizes those commutative
automata with longest possible synchronizing words.

\subsection{Commutativity of Rees Quotients}
\label{prelim:classes:comm:subsec:rees}

\begin{proposition} \label{prelim:subclasses::com:upper-sets:prop} 
	Let  $\semiautomat{A}{Q}{\Sigma}{\delta}$ and $S \subseteq Q$ be an upper set.
	Then the following are equivalent.
	\begin{enumerate}
		\item \label{prelim:subclasses::com:upper-sets:prop:1} $\mathcal A / {\sim_S}$ is commutative.
		
		\item \label{prelim:subclasses::com:upper-sets:prop:2} $\delta(q, ab) = \delta(q, ba)$ for each $q \in Q \setminus S$ and $a,b \in \Sigma$ with $\delta(q, ab) \in Q \setminus S$.
		
		\item \label{prelim:subclasses::com:upper-sets:prop:3} $\delta(q, ab) \ne \delta(q, ba)$
		implies
		$
		\{ \delta(q, ab), \delta(q, ba) \} \subseteq S
		$
		for all $q \in Q$ and $a,b \in \Sigma$.
	\end{enumerate}
\end{proposition}
\begin{proof}
	$\ref{prelim:subclasses::com:upper-sets:prop:1} \Rightarrow \ref{prelim:subclasses::com:upper-sets:prop:2}$: This directly follows from \cref{prelim:rem:Rees}.
	
	$\ref{prelim:subclasses::com:upper-sets:prop:2} \Rightarrow \ref{prelim:subclasses::com:upper-sets:prop:3}$: By Item~\ref{prelim:subclasses::com:upper-sets:prop:1},
	for $\{ \delta(q, ab), \delta(q, ba) \} \not\subseteq S$,
	we have  $\delta(q, ab) = \delta(q, ba)$,
	which gives the claim.
	
	$\ref{prelim:subclasses::com:upper-sets:prop:3} \Rightarrow \ref{prelim:subclasses::com:upper-sets:prop:1}$: Since~$S$ is a trap state in $\mathcal A / {\sim_S}$,
	from Item~\ref{prelim:subclasses::com:upper-sets:prop:3} and \cref{prelim:rem:Rees} we obtain $\delta_{\sim_S}(\pi_{\sim_S}(p), ab) = \delta_{\sim_S}(\pi_{\sim_S}(q), ba)$
	for all $q \in Q$ and $a,b \in \Sigma$, i.e., $\mathcal A / {\sim_S}$ is commutative.	
\end{proof}

\begin{corollary} \label{prelim:classes:comm:prop:perm-reset-comm}
	Let  $\semiautomat{A}{Q}{\Sigma}{\delta}$ and $S \subseteq Q$ be an upper set.
	Then $\mathcal A / {\sim_S}$ is a commutative permutation-reset automaton
	if and only if
	$\restrictAutGamS{\mathcal A}{\Gamma}{Q \setminus S}$
	for  $\Gamma = \{ a \in \Sigma \mid \delta(Q \setminus S, a) = Q \setminus S \}$
	is commutative and $\delta(Q \setminus S, a) \subseteq S$
	for $a \in \Sigma \setminus \Gamma$.	
\end{corollary}
\begin{proof}
	First, suppose $\mathcal A / {\sim_S}$ is a commutative permutation-reset automaton.
	Then clearly $\restrictAutGamS{\mathcal A}{\Gamma}{Q \setminus S}$
	is commutative. By \cref{prelim:perm-reset-Rees-factor} we have 
	$\delta(Q \setminus S, a) \subseteq S$ for $a \in \Sigma \setminus \Gamma$.
	Conversely, since each letter in $\Gamma$ permutes $Q \setminus S$,
	$\restrictAutGamS{\mathcal A}{\Gamma}{Q \setminus S}$ is a permutation automaton.
	Thus, with \cref{prelim:perm-reset-Rees-factor},
	we find that $\mathcal A / {\sim_S}$ is a permutation-reset automaton.
	Next, let $q \in Q \setminus S$ and $a,b \in \Sigma$.
	For $a,b \in \Gamma$, we have $\delta(q, ab) = \delta(q, ba)$
	since 	$\restrictAutGamS{\mathcal A}{\Gamma}{Q \setminus S}$
	is commutative.
	Next, without loss of generality, suppose $a \notin \Gamma$.
	Then $\delta(Q, a) \subseteq S$
	and this gives $\{ \delta(q, ab), \delta(q, ba) \}\subseteq S$.
	Thus, by \cref{prelim:subclasses::com:upper-sets:prop}, 
	$\mathcal A / {\sim_S}$ is commutative.	
\end{proof}

\begin{corollary} \label{cor:SCC_comm_permutation_automaton}
	Let  $\semiautomat{A}{Q}{\Sigma}{\delta}$ and $S \subseteq Q$ be an upper set
	such that  $\mathcal A / {\sim_S}$ is commutative and $Q \setminus S$ is strongly connected.
	Then $\mathcal A / {\sim_S}$ is a permutation-reset automaton (resp.~permutation automaton for $S = \emptyset$).	
\end{corollary}
\begin{proof}
	Since~$S$ is an upper set, $Q \setminus S$ is a strongly connected component.
	By \cref{prelim:automata:Rees:prop:correspondence},
	$\pi_{\sim_S}(Q \setminus S)$ is a strongly connected component in $\mathcal A / {\sim_S}$.
	Thus, $\pi_{\sim_S}(Q \setminus S)$ is a block by 
	\cref{prelim:automata:cor:SCC-disjoint-or-permutates}.
	Now, a set $R \subseteq Q \setminus S$ is a block in $\mathcal A$ if and only if $\pi_{\sim_S}(R)$
	is a block in $\mathcal A / {\sim_S}$. 
	Thus, $Q \setminus S$ is also a block in $\mathcal A$. 
	For $S = \emptyset$, this yields that~$\mathcal A / {\sim_S}$
	is a permutation automaton
	and, by \cref{prelim:perm-reset-Rees-factor}, we find that $\mathcal A / {\sim_S}$
	is a permutation-reset automaton.	
\end{proof}

\subsection{Strongly Connected Automata}
\label{prelim:classes:comm:subsec:SCC}

\begin{proposition} \label{prelim:classes:comm:cor:collections-of-facts:7} 
	A commutative automaton that is strongly connected (resp.~strongly trap-connected)
	is a permutation automaton (resp.~permutation-reset automaton with a trap state). 
\end{proposition}
\begin{proof} 
	This follows from \cref{cor:SCC_comm_permutation_automaton} applied with 
	$S = \emptyset$ or with $S = \{t\}$ for a trap state.
\end{proof} 

\begin{corollary}
	A strongly connected commutative automaton with at least two states is not synchronizing.	
\end{corollary}
\begin{proof}
	This follows from \cref{prelim:classes:comm:cor:collections-of-facts:7} 
	since a synchronizing permutation automaton has a single state.
\end{proof}

\begin{proposition} \label{prelim:classes:com:prop:endos}
	In a cyclic commutative automaton $\semiautomat{A}{Q}{\Sigma}{\delta}$,
	the endomorphism semigroup of $\mathcal A$ equals
	$\TM{\mathcal A}$.
\end{proposition}
\begin{proof}
	In a commutative automaton, every element of its transformation monoid
	is an endomorphism.
	By assumption, we have $Q = \upperset{\{q\}}$ for some $q \in Q$.
	Now, let $\varphi \colon Q \to Q$ be an endomorphism.
	Then $\varphi(q) = \delta(q, u)$ for some $u \in \Sigma^*$.
	Consider some state $p \in Q$. We have $p = \delta(q, v)$
	for some $v \in \Sigma^*$.
	This gives $\varphi(p) = \varphi(\delta(q, v)) = \delta(\varphi(q), v) = \delta(q, uv) = \delta(\delta(q, v), u) = \delta(p, u)$
	and so $\varphi = \fword{\delta}{u}$.
\end{proof}

\begin{corollary} \label{prelim:classes:comm:prop:permSCC-regular}
	Let $\semiautomat{A}{Q}{\Sigma}{\delta}$ be a strongly connected commutative automaton.
	Then $\TM{\mathcal A}$ is a regular permutation group and thus has size~$|Q|$.
\end{corollary}
\begin{proof}
	Since~$\mathcal A$ is strongly connected,
	its endomorphism semigroup
	equals it automorphism group by	\cref{prelim:automata:morphism:prop:endo-SCC-is-iso}.
	By \cref{prelim:automata:morphism:prop:endo-SCC-is-iso-b}
	the automorphism group is semiregular.
	By \cref{prelim:classes:com:prop:endos} 
	it equals $\TM{\mathcal A}$.
	Since $\mathcal A$ is strongly connected, $\TM{\mathcal A}$
	is transitive. Thus, $\TM{\mathcal A}$ is regular. 
\end{proof}

Next, we give examples of commutative automata 
where the transformation monoid is properly contained in its endomorphism monoid. 

\begin{example}
	Let $\semiautomat{A}{Q}{\{a_1,a_2\}}{\delta}$ 
	with $Q = \{ q_{1,1}, q_{1,2}, q_{2,1}, q_{2,2} \}$
	where, for $i \in \{1,2\}$, we
	set $\delta(q_{i,1}, a_i) = \delta(q_{i,2}, a_i) = q_{i,2}$
	and the remaining transitions are self-loops.
	Then $\mathcal A$ is commutative and 
	$\varphi \colon Q \to Q$ with $\varphi(q_{i,1}) = \varphi(q_{i,2}) = q_{i,2}$
	for $i \in \{1,2\}$ is a morphism that is not contained
	in $\TM{\mathcal A}$.	
\end{example}

\begin{example}
	Let $\semiautomat{A}{\{ q_{i,j} \mid i \in \{0,1\},\, j \in [n] \}}{\{a\}}{\delta}$
	be the semiautomaton with
	\[
	\delta(q_{i,j}, a) = q_{i,(j+1)\bmod n}
	\]
	for $i \in \{0,1\}$ and $j \in [n]$.
	For every $0 \le s,t < n$, the map $\varphi$ defined by
	\[
	\varphi(q_{0,j}) = q_{1,(j+s)\bmod n},
	\qquad
	\varphi(q_{1,j}) = q_{0,(j+t)\bmod n}
	\]
	is an endomorphism since
	\[
	\delta(\varphi(q_{0,j}),a)
	= q_{1,(j+s+1)\bmod n}
	= \varphi(q_{0,(j+1)\bmod n})
	= \varphi(\delta(q_{0,j},a)),
	\]
	and similarly for the states $q_{1,j}$.
	None of these endomorphisms is induced by a word and thus
	none belongs to the transformation monoid.
\end{example}

\label{def:grid-aut}
Next, we determine the structure of binary commutative strongly connected automata.
Let $n,m \ge 1$ and $r \in [n]$.
Define automata $\semiautomat{B_{n,m, r}}{[n] \times [m]}{\{a,b\}}{\mu_{n,m,r}}$
where $\mu_{n,m, r}((i,j), a) = ( (i+1) \bmod{n}, j )$
and 
\[
\mu_{n,m, r}((i,j), b)
= \begin{cases}
	(i, j + 1) & \text{if } 0 \le j \le m - 2, \\
	((i + r) \bmod n, 0) & \text{if } j = m - 1. 		
\end{cases}
\]	
for $(i,j) \in [n] \times [m]$. These automata are commutative. See \cref{prelim:fig:Bnmr} for some examples.

\begin{remark}
	For $n,m \ge 2$, the automaton $\mathcal B_{n,m,1}$ looks like
	a grid where each component counts modulo~$n$
	and~$m$ the number of~$a$'s and~$b$'s, respectively,
	and each state is uniquely determined by these numbers.
\end{remark}

\begin{remark}	
	In $\mathcal B_{n,1,r}$ the letter $a$ induces a single $n$-cycle and we have
	\[
	\mu_{n,1,r}((i, 1), b) = \mu_{n,1,r}((i, 1), a^r), \quad i \in [n],
	\]
	i.e., the letter $b$ ``wraps'' around this cycle by skipping $r$ steps on the cycle.
\end{remark}

\begin{figure}[htbp]
	\centering
	\begin{minipage}[b]{0.3\textwidth}
		\centering
		
		\begin{tikzpicture}[>=Stealth, on grid, auto]
			
			\tikzset{
				state/.style={
					circle,
					draw,
					minimum size=6mm,
					inner sep=1pt,
					font=\small
				}
			}
			
			\node[state] (q0) at (90:5em) {$(0,0)$};
			\node[state] (q1) at (30:5em) {$(1,0)$};
			\node[state] (q2) at (-30:5em) {$(2,0)$};
			\node[state] (q3) at (-90:5em) {$(3,0)$};
			\node[state] (q4) at (-150:5em) {$(4,0)$};
			\node[state] (q5) at (150:5em) {$(5,0)$};
			
			\path
			(q0) edge[dashed, bend left=12] (q1)
			(q1) edge[dashed, bend left=12] (q2)
			(q2) edge[dashed, bend left=12] (q3)
			(q3) edge[dashed, bend left=12] (q4)
			(q4) edge[dashed, bend left=12] (q5)
			(q5) edge[dashed, bend left=12] (q0);
			
			\path
			(q0) edge (q2)
			(q1) edge (q3)
			(q2) edge (q4)
			(q3) edge (q5)
			(q4) edge (q0)
			(q5) edge (q1);	
		\end{tikzpicture}
		\caption*{$\mathcal B_{6,1,2}$}
	\end{minipage}%
	\hfill
	\begin{minipage}[b]{0.3\textwidth}
		\centering
		\begin{tikzpicture}[>=Stealth]
			
			\tikzset{
				state/.style={
					circle,
					draw,
					minimum size=6mm,
					inner sep=1pt,
					font=\small
				}
			}
			
			\matrix (m) [matrix of nodes,
			nodes={state},
			row sep=2.2em,
			column sep=4.5em
			]{
				$(0,0)$ & $(0,1)$ \\
				$(1,0)$ & $(1,1)$ \\
				$(2,0)$ & $(2,1)$ \\
				$(3,0)$ & $(3,1)$ \\
			};
			
			\path
			(m-1-1) edge[dashed] (m-2-1)
			(m-2-1) edge[dashed] (m-3-1)
			(m-3-1) edge[dashed] (m-4-1)
			(m-4-1) edge[dashed,bend left=25] (m-1-1)
			
			(m-1-2) edge[dashed] (m-2-2)
			(m-2-2) edge[dashed] (m-3-2)
			(m-3-2) edge[dashed] (m-4-2)
			(m-4-2) edge[dashed,bend right=25] (m-1-2);
			
			\path
			(m-1-1) edge[bend left=25] (m-1-2)
			(m-2-1) edge[bend left=25] (m-2-2)
			(m-3-1) edge[bend right=25] (m-3-2)
			(m-4-1) edge[bend right=25] (m-4-2)
			
			(m-1-2) edge[bend left=10] (m-3-1)
			(m-2-2) edge[bend right=10] (m-4-1)
			(m-3-2) edge[bend left=10] (m-1-1)
			(m-4-2) edge[bend right=10] (m-2-1);
			
		\end{tikzpicture}	
		\caption*{$\mathcal B_{4,2,2}$}
	\end{minipage}%
	\hfill
	\begin{minipage}[b]{0.3\textwidth}
		\centering
		\begin{tikzpicture}[>=Stealth, auto]
			
			\tikzset{
				state/.style={
					circle,
					draw,
					minimum size=6mm,
					inner sep=1pt,
					font=\small
				}
			}
			
			\node[state] (q0) at (90:5em)  {(0,0)};
			\node[state] (q1) at (30:5em)  {(0,1)};
			\node[state] (q2) at (-30:5em) {(2,0)};
			\node[state] (q3) at (-90:5em) {(2,1)};
			\node[state] (q4) at (-150:5em){(1,0)};
			\node[state] (q5) at (150:5em) {(1,1)};
			
			\draw (q0) to[bend left=15] (q1);
			\draw (q1) to[bend left=15] (q2);
			\draw (q2) to[bend left=15] (q3);
			\draw (q3) to[bend left=15] (q4);
			\draw (q4) to[bend left=15] (q5);
			\draw (q5) to[bend left=15] (q0);
			
			\draw[dashed] (q0) to (q4);
			\draw[dashed] (q1) to (q5);
			\draw[dashed] (q2) to (q0);
			\draw[dashed] (q3) to (q1);
			\draw[dashed] (q4) to (q2);
			\draw[dashed] (q5) to (q3);	
		\end{tikzpicture}
		\caption*{$\mathcal B_{3,2,2}$}
	\end{minipage}
	\caption[Example automata $\mathcal B_{n,m,r}$]{Automata $\mathcal B_{n,m,r}$ with $a$-transitions drawn dashed and $b$-transitions solid.}
	\label{prelim:fig:Bnmr}
\end{figure}

\begin{remark}	
	Suppose $r$ is coprime to $n$, and let $q=(0,0)$.
	Then there exists an integer $l$ with $1 \le l < n$ such that
	$rl \equiv 1 \pmod n$, i.e., $l$ is the multiplicative inverse of $r$
	modulo $n$.
	We have the following properties:
	
	\begin{enumerate}
		\item The states
		\[
		q,\delta(q,b),\ldots,\delta(q,b^{nm})
		\]
		are pairwise distinct and form a single cycle.
		
		\item Since $\delta(q,b^m)=\delta(q,a^r)$, we obtain
		\[
		\delta(q,a)
		= \delta(q,a^{rl})
		= \delta(q,b^{ml}).
		\]
		Consequently, we have
		\[
		\delta(q,b^i a)
		= \delta(q,a b^i)
		= \delta(q,b^{i+ml})
		\]
		for all $i \ge 0$. In other words, the letter $a$ acts on the $b$-cycle
		as a shift by $ml$ positions.
	\end{enumerate}
	
	By these properties, the automaton has the same structure as
	$\mathcal B_{nm,1,ml}$, with the roles of $a$ and $b$ interchanged.	
\end{remark}

\begin{proposition} \label{prelim:classes:com:prop:comm-perm-classifiy0}
	Let $\mathcal A = (Q, \{a,b\}, \delta)$ be a commutative and strongly connected.
	Then $\mathcal A$ is isomorphic to 
	$\mathcal B_{n,m,r}$ for some $n,m \ge 1$
	and $r \in [n]$.
\end{proposition}
\begin{proof}
	By \cref{prelim:classes:comm:cor:collections-of-facts:7}, $\mathcal A$
	is a permutation automaton.
	Let $q \in Q$
	and $n \ge 0$ be minimal
	such that $\delta(q, a^n) = q$.
	Then the set 
	\begin{equation}\label{prelim:classes:com:eqn:Bnm}
	S = \{q, \delta(q, a), \ldots, \delta(q, a^{n-1})\}
	\end{equation}
	forms a strongly connected component in $\mathcal A_{\{a\}}$.
	The letter~$b$ acts as an isomorphism on this automaton.
	By \cref{prelim:automata:morphism:prop:injective:SCC-to-SCC} from \cref{prelim:automata:morphism:prop:injective}
	it permutes the strongly connected components of $\mathcal A_{\{a\}}$.
	Thus, we have $\delta(S, b^m) = S$ for some minimal $m \ge 1$.  
        
    Next, we show that
    \begin{equation}\label{prelim:classes:com:eqn:Bnm-S}
    Q = S \cup \delta(S, b) \cup \cdots \cup \delta(S, b^{m-1}).
    \end{equation}
    Let $p \in Q$.
    Since~$\mathcal A$ is commutative and strongly connected,
	every state is reachable from~$q$ by some word from~$a^* b^*$.
	Thus $p = \delta(q, a^i b^j)$ for some $i,j \ge 0$
	where these numbers are picked such that~$j$ is minimal.
	Assume $j > m$.
	We have $\delta(q, a^i b^m) \in S$.
	By \cref{prelim:classes:com:eqn:Bnm}, this yields $\delta(q, a^i b^m) = \delta(q, a^l)$ for some $0 \le l < n$.
	Hence $p = \delta(q, a^l b^{j - m})$, contradicting the minimal choice of~$j$.
	Thus, we have $j < m$. So, $p \in \delta(S, b^j)$ with $0 \le j < m$.

	Now, consider $\varphi \colon [n] \times [m] \to Q$
	given by $\varphi(i,j) = \delta(q, a^i b^j)$
	for $(i,j) \in [n] \times [m]$. By \cref{prelim:classes:com:eqn:Bnm,prelim:classes:com:eqn:Bnm-S}
	this map is surjective.	Next, let
	$\varphi(i,j) = \varphi(i', j')$ for $(i,j), (i',j') \in [n] \times [m]$.
	Then $\varphi(i',j') \in \delta(S, b^j)$.
	Since the sets $S, \delta(S, b), \ldots, \delta(S, b^{m-1})$
	are pairwise disjoint,
	we must have $j' = j$.
	Since~$b^j$ induces a permutation on the state set,
	we get $\delta(q, a^i) = \delta(q, a^{i'})$.
	Since $i,i' \in [n]$ and $q, \delta(q,a), \ldots, \delta(q,a^{n-1})$
	are pairwise distinct, we find $i = i'$.
	Thus, $\varphi$ is also injective.
	
	Since $\delta(S, b^m) = S$ and $q \in S$, 
	by \cref{prelim:classes:com:eqn:Bnm} there exists $r \in [n]$ such that
	\begin{equation} \label{prelim:classes:com:eqn:Bnm-r}
	\delta(q, b^m) = \delta(q, a^r).
	\end{equation}
	Next, we show that $\varphi$ is a morphism
	of $\mathcal A$ and $\mathcal B_{n,m,r} = ([n] \times [m], \{a,b\}, \mu_{n,m,r})$.	
	Let $(i,j) \in [n] \times [m]$.
	We have 
	\begin{align*}
	\delta(\varphi(i,j), a)
	& = \delta(\delta(q, a^i b^j), a)  
	  = \delta(q, a^{i+1}b) \tag{Def.~$\varphi$ and $\delta$} \\
	& = \varphi((i+1)\bmod n, j) \tag{Eqn.~\eqref{prelim:classes:com:eqn:Bnm} and Def.~$\varphi$} \\ 
	& = \varphi(\mu_{n,m,r}((i,j), a)). \tag{Def.~$\mu_{n,m,r}$}
	\end{align*}
	For $j \le m - 2$, we similarly get $\delta(\varphi(i,j), b) = \varphi(\mu_{n,m,r}((i,j), b))$.
	Lastly, we have
	\begin{align*}
	\delta(\varphi(i,m-1), b) 
    & = \delta(q, a^i b^m) = \delta(q, b^m a^i) \tag{Def.~$\varphi$ and $\mathcal A$ comm.} \\ 
	& = \delta(q, a^{i + r}) \tag{Eqn.~\eqref{prelim:classes:com:eqn:Bnm-r}} \\
	& = \varphi((i+r) \bmod{n}, 0)  \tag{Eqn.~\eqref{prelim:classes:com:eqn:Bnm} and Def.~$\varphi$} \\
	& = \varphi(\mu_{n,m,r}((i,m-1), b)). \tag{Def.~$\mu_{n,m,r}$}	
	\end{align*}
	Thus $\varphi$ is an isomorphism between $\mathcal A$ and $\mathcal B_{n,m,r}$.
\end{proof}

\begin{corollary}
	\label{prelim:classes:com:prop:comm-perm-classifiy}
	Let $\mathcal A = (Q, \{a,b\}, \delta)$ be a commutative permutation automaton
	and $S \subseteq Q$ be a strongly connected component.
	Then~$S$ is a sink component
	and $\mathcal A_{|S} \cong \mathcal B_{n,m,r}$
	for some $n,m \ge 1$ and $r \in [n]$.
\end{corollary}
\begin{proof}
	By \cref{prelim:subclasses:perm:reach:cor}, $S$ is a sink component.
	By \cref{prelim:classes:com:prop:comm-perm-classifiy0}, $\mathcal A_{|S}$
	isomorphic to $\mathcal B_{n,m,r}$ for some $n,m \ge 1$
	and $r \in [n]$.	
\end{proof}

\begin{remark} \label{prelim:classes:com:rem:comm-perm-classify}
	Let $\varphi \colon [n] \times [m] \to [n] \times [m]$
	be an automorphism of $\mathcal B_{n,m,r}$.
	Then $\varphi((0,0)) = (s,t) = \mu_{n,m,r}((0,0), a^sb^t)$
	for some $(s,t) \in [n] \times [m]$
	and this value determines $\varphi$ uniquely
	by \cref{prelim:automata:morphism:prop:principal-upper-set}.
	Thus, the automorphism group of $\mathcal B_{n,m,r}$
	consists of all mappings $(i,j) \mapsto \mu((i,j), a^s b^t)$
	for $(s,t) \in [n] \times [m]$.
	We see that this group equals $\TM{\mathcal B_{n,m,r}}$
	and is regular.
	With \cref{prelim:classes:com:prop:comm-perm-classifiy0} this gives another proof of
	\cref{prelim:classes:comm:prop:permSCC-regular} 
\end{remark}

\subsection{Subset Mapping Properties}
\label{prelim:classes:comm:subsec:subset}

We begin with some statements on synchronizing commutative automata.

\begin{proposition}[{Fernau and Hoffmann~\cite[Lem.~20]{FernauHoffmann19}}] \label{prelim:classes:comm:cor:collections-of-facts:2}
 In a commutative automaton, a synchronizing state is a trap state.
\end{proposition} 
\begin{proof}
	This follows from \cref{constr_sync:prop:sink-blocks-sync-trap} since the sink components are blocks.	
\end{proof}

\begin{corollary} \label{prelim:classes:comm:cor:sync-trap}
	Let $\mathcal A = (Q, \Sigma, \delta)$ be a commutative permutation-reset automaton
	that is not a permutation automaton.
	Then there exists a trap state~$t \in Q$
	such that for each $a \in \Sigma$ that is not a permutational letter
	we have $\delta(Q, a) = \{t\}$.	
\end{corollary}
\begin{proof}
	By assumption, $\mathcal A$ has a letter mapping all states
	to the same state~$t$.
	By \cref{prelim:classes:comm:cor:collections-of-facts:2}, $t$ is a trap state.
	Let $a \in \Sigma$ be some state that is not a permutational letter.
	Then $|\delta(Q, a)| = 1$ and, since~$t$ is a trap state,
	$t \in \delta(Q, a)$. This gives $\delta(Q, a) = \{t\}$.
\end{proof}

\begin{proposition}  \label{prelim:classes:comm:cor:collections-of-facts:6}
	$\mathcal A = (Q, \Sigma, \delta)$ with $|Q| \ge 2$ is a synchronizing commutative permutation-reset
	automaton if and only if there exists a non-permutational letter, a trap state $t \in Q$,
	and $S = Q \setminus \{t\}$ is a block
	such that $\delta(q, ab) = \delta(q, ba)$ for all $q \in S$ and $a,b\in \Sigma$
	with $\delta(q,ab) \in S$.
\end{proposition}
\begin{proof}	
	First, suppose~$\mathcal A$
	is a synchronizing commutative permutation-reset automaton.
	By \cref{prelim:classes:comm:cor:sync-trap}
	we have a trap state $t$, a non-permutational letter, and for the upper set $S = \{t\}$
	we have $\delta(Q, a) = Q$ or $\delta(Q, a) = \{t\}$
	for each $a \in \Sigma$. Thus, $S$ is a block.
	Since~$\mathcal A$ is commutative, we have 
	$\delta(q, ab) = \delta(q, ba)$ for all $q \in S$ and $a,b\in \Sigma$.
		
	Conversely, by considering $\mathcal A / {\sim_{\{t\}}}$, which is isomorphic to~$\mathcal A$,
	from \cref{prelim:perm-reset-Rees-factor,prelim:subclasses::com:upper-sets:prop}
	we find that it is a commutative permutation-reset automaton,
	which is synchronizing since it has a non-permutational letter
	mapping everything to~$t$. 
\end{proof}

\begin{proposition} \label{prelim:classes:comm:cor:collections-of-facts:5} 
	A two-state automaton is commutative and synchronizing 
	if and only if it has a trap state~$t$,
	some non-identity letter, and every letter acts as the identity
	or maps both states to~$t$.
\end{proposition}
\begin{proof}
	A synchronizing and commutative automaton
	has a trap state~$t$ by \cref{prelim:classes:comm:cor:collections-of-facts:2}.
	Thus, each letter either fixes both states or maps both states to~$t$.
	Conversely, an automaton of this form is easily seen to be synchronizing and commutative.	
\end{proof}

\begin{proposition} \label{prelim:classes:comm:lem:words-SIS-SFS}
	Let $\semiautomat{A}{Q}{\Sigma}{\delta}$
	be a commutative $n$-state automaton and $u \in \Sigma^*$.
	\begin{enumerate}
		\item \label{prelim:classes:comm:lem:words-SIS-SFS:1} We have $\delta(Q, u) = \delta(Q, v)$
		for all $v \in \Sigma^*$
		with $|v|_a = \min\{ n - 1, |u|_a \}$
		for $a \in \Sigma$.
		
		\item \label{prelim:classes:comm:lem:words-SIS-SFS:2} For $S \subseteq Q$, $|\delta(S, u)| = |\delta(S, v)|$ for all $v \in \Sigma^*$
		with $|v|_a = \min\{ n - 1, |u|_a \}$ for $a \in \Sigma$.
	\end{enumerate}
\end{proposition}
\begin{proof}
	\ref{prelim:classes:comm:lem:words-SIS-SFS:1} 
	Let  $\Gamma \subseteq \Sigma$ be a largest alphabet
	such that, for all $v \in \Sigma^*$
	with $|v|_a = \min\{ n - 1, |u|_a\}$
	for $a \in \Gamma$
	and $|v|_a = |u|_a$
	for $a \in \Sigma \setminus \Gamma$,
	we have $\delta(Q, u) = \delta(Q, v)$.
	Since, by commutativity, $\Gamma = \emptyset$ 
	fulfills this condition, such a set~$\Gamma$
	exists.	
	
	Assume there exist $a \in \Sigma \setminus \Gamma$.
	Let $w \in \Sigma^*$
	with $|w|_b = \min\{n-1, |u|_b\}$ for $b \in \Gamma$,
	$|w|_b = |u|_b$ for $b \in \Sigma \setminus \Gamma$
	with $b \ne a$,
	and $|w|_a = 0$.
	By choice of~$\Gamma$,
	we have $\delta(Q, u) = \delta(Q, a^{|u|_a} w)$.
	By \cref{prelim:automata:cor:power:P} 
	we find $\delta(Q, a^{|u|_a} w) = \delta(Q, a^{\min\{n - 1, |u|_a\}} w)$.
	Set $w' = a^{\min\{n - 1, |u|_a\}} w$.
	By commutativity,
	we have $\delta(Q, w') = \delta(Q, v)$
	for all $v \in \Sigma^*$
	with $|v|_b = | w' |_b$
	for $b \in \Sigma$.
	Since $|w|_a = \min\{n-1,|u|_a\}$, $|w|_b = \min\{n-1,|u|_b\}$ for $b \in \Gamma$,
	and $|w|_b = |u|_b$ otherwise,
	and $\delta(Q, u) = \delta(Q, v)$,
	this contradicts the maximal choice of~$\Gamma$
	since the condition by which we selected it 
	is also fulfilled for $\Gamma \cup \{a\}$.
	Thus, we must have $\Sigma = \Gamma$, which gives the claim.
		
	\ref{prelim:classes:comm:lem:words-SIS-SFS:2}: Let $S \subseteq Q$. We show, by induction,
	that for each $\Gamma \subseteq \Sigma$ 
	we have $|\delta(S, u)| = |\delta(S, v)|$ for all $v \in \Sigma^*$
	with $|v|_a = \min\{n - 1, |u|_a\}$ for $a \in \Gamma$
	and $|v|_a = |u|_a$ for $a \in \Sigma \setminus \Gamma$.
	Then picking $\Gamma = \Sigma$ gives the statement we want to prove.
	The claim clearly holds for $\Gamma = \emptyset$. 
	
	Next, consider $a \in \Sigma \setminus \Gamma$.
	Let $v \in \Sigma^*$ with $|v|_b = \min\{ n - 1, |u|_b \}$
	for $b \in \Gamma \cup \{a\}$ and $|v|_b = |u|_b$ for $b \in \Sigma \setminus \Gamma$ with $b \ne a$.
	Choose $w \in \Sigma^*$ 
	such that $|w|_b = \min\{n - 1, |u|_b\}$ for $b \in \Gamma$,
	$|w|_b = |u|_b$ for $b \in \Sigma \setminus \Gamma$ with $b \ne a$
	and $|w|_a = 0$.
	Then, by induction hypothesis, we have
	$
	|\delta(S, u)| = |\delta(S, w a^{|u|_a})|.
	$
	From \cref{prelim:automata:power:prop:S:3} of \cref{prelim:word_in_PA},
	for $v' = w a^{\min\{n-1,|u|_a\}}$ we get
	\[
	|\delta(S, w a^{|u|_a})| = |\delta(S, v')|.
	\]
	We have $|v|_b = |v'|_b$ for all $b \in \Sigma$. 
	Thus, by commutativity, $\delta(S, v) = \delta(S, v')$
	and so $|\delta(S, u)| = |\delta(S, v)|$.  This concludes
	the induction step.
\end{proof}

\begin{corollary}
	In a commutative $n$-state automaton, a shortest word reaching
	a given subset 
	has length at most $|\Sigma|(n-1)$.	The same bound holds
	for shortest words synchronizing a given subset in commutative automata.
\end{corollary}

Note that \cref{prelim:classes:block:prop:sync-into-bound} gives 
a better bound than $|\Sigma|(n-1)$ in commutative automata 
for words mapping the whole state set into a given subset.


\begin{corollary} \label{prelim:classes:comm:cor:words-SIS-SFS}
	Let $\Sigma = \{a_1, \ldots, a_k\}$.
	Suppose $\mathcal A = (Q, \Sigma,\delta)$ is a commutative $n$-state automaton
	and $u \in \Sigma^*$.
	Then we have
	\begin{enumerate}
	\item $\delta(Q, a_1^{n-1} \cdots a_k^{n-1}) \subseteq \delta(Q, u)$ and 
		
	\item $|\delta(S, a_1^{n-1} \cdots a_k^{n-1})| \le |\delta(S, u)|$ for $S \subseteq Q$.
	\end{enumerate}
\end{corollary}
\begin{proof}
	Set $v = a_1^{\min\{n-1, |u|_{a_1}\}} \cdots a_k^{\min\{n-1, |u|_{a_k}\}}$
	and $w = a_1^{n - 1 - |v|_{a_1}} \cdots a_k^{n-1 - |v|_{a_k}}$.
	By \cref{prelim:classes:comm:lem:words-SIS-SFS} 
	we have $\delta(Q, u) = \delta(Q, v)$
	and $|\delta(S, u)| = |\delta(S, v)|$.
	Since $\delta(Q, w) \subseteq Q$
	we have $\delta(Q, wv) \subseteq \delta(Q, v)$
	and $|\delta(S, vw)| \le |\delta(S, v)|$.
	By commutativity,
	we find $\delta(Q, wv) = \delta(Q, a_1^{n-1} \cdots a_k^{n-1})$
	and $\delta(S, vw) = \delta(S, a_1^{n-1} \cdots a_k^{n-1})$,
	and this yields the claim.	
\end{proof}

By \cref{prelim:classes:comm:cor:words-SIS-SFS},
in a commutative $n$-state automaton we have $\delta(Q, u) \subseteq S$ (resp.~$|\delta(S, u)| = 1$)
for some $u \in \Sigma^*$
if and only if $\delta(Q, a_1^{n-1} \cdots a_k^{n-1}) \subseteq S$ 
(resp.~$|\delta(S, a_1^{n-1} \cdots a_k^{n-1})| = 1$).

Because it easily follows from other results,
we state a tight bound for shortest synchronizing words 
that was first obtained by Rystsov~\cite[Thm.~3.5]{Rys96} and~\cite[Thm.~3.2]{Rystsov97} 
(see also~\cite[Cor.~19]{FernauHoffmann19} for another proof).

\begin{proposition} \label{prelim:classes:comm:prop:set-sync-words}
	Let $\mathcal A = (Q, \Sigma,\delta)$ be commutative and $n = |Q|$.
	\begin{enumerate}
		\item \label{prelim:classes:comm:prop:set-sync-words:1} $\Syn(\mathcal A)$ is commutative.
		
		\item \label{prelim:classes:comm:prop:set-sync-words:2} $uwv \in \Syn(\mathcal A)$ for all $u,v,w \in \Sigma^*$ with $uv \in \Syn(\mathcal A)$.
		
		\item \label{prelim:classes:comm:prop:set-sync-words:3} (Rystsov~\cite{Rys96,Rystsov97}) A shortest word in $\Syn(\mathcal A)$ has length at most $n - 1$ and this bound is tight.		
	\end{enumerate}
\end{proposition}
\begin{proof}
	$\ref{prelim:classes:comm:prop:set-sync-words:1}$: By commutativity, $\delta(Q, uabv) = \{q\} = \delta(Q, ubav)$
	for all $u,v \in \Sigma^*$ and $a,b \in \Sigma$,
	which yields that $\Syn(\mathcal A)$ is commutative.
	
	$\ref{prelim:classes:comm:prop:set-sync-words:2}$: Since~$\Syn(\mathcal A)$ is an ideal by \cref{prelim:thm:sync_words_is_ideal},
	we have $uvw \in \Syn(\mathcal A)$. By Item~\ref{prelim:classes:comm:prop:set-sync-words:1},
	we get $uwv \in \Syn(\mathcal A)$.
	
	$\ref{prelim:classes:comm:prop:set-sync-words:3}$: Since in a commutative automaton
	the strongly connected components are blocks,
	this follows from \cref{prelim:thm:SCCblocks-rank_r} with $r = 1$,
	or from  \cref{prelim:classes:block:prop:sync-into-bound} applied $S = \{t\}$
	where $t$ is the synchronizing state, which must be a trap state
	so that~$S$ is an upper set.
\end{proof}

Next, in \cref{syncmax:classes:comm:charsync}, 
we show that weak connectivity and existence of a trap state is sufficient
for synchronizability in commutative automata. This follows from a couple of statements
we give next.

\begin{lemma} \label{lem:comm:inv-images-letter-upper-set}
	Let $\semiautomat{A}{Q}{\Sigma}{\delta}$ be a commutative automaton
	with sink component $S \subseteq Q$ and $\Gamma \subseteq \Sigma$.
	Then the set $T = \{ q \in Q \mid \exists u \in \Gamma^* \colon \delta(q, u) \in S \}$
	is an upper set
	such that $S$ is the only sink component having non-empty intersection with~$T$.
\end{lemma}
\begin{proof}
	Let $q \in T$ and $b \in \Sigma$.
	We have $\delta(q, u) \in S$ for some $u \in \Gamma^*$.
	By commutativity, we find $\delta(q, b u) = \delta(q, u b)$.
	Since $\delta(q, u) \in S$ and~$S$ is a sink component
	we have $\delta(q, ub) \in S$.	
	Thus, $\delta(q, b u) \in S$ and so $\delta(q, b) \in T$. Hence, $T$ is an upper set.
	
	Next, let $R \subseteq Q$ be a sink component with $R \cap T \ne \emptyset$.
	By the definition of~$T$, 
	we have $\delta(R, u) \cap S \ne \emptyset$ for some $u \in \Gamma^*$.
	Since~$R$ is a sink component, we have $\delta(R, u) \subseteq R$.
	Thus $R \cap S \ne \emptyset$.
	Since distinct sink components are disjoint, this gives $R = S$.	
\end{proof}

\begin{corollary} \label{cor:comm:inv-images-letter-upper-set-T-weakly}
	Let $\semiautomat{A}{Q}{\Sigma}{\delta}$ be commutative
	with sink component $S \subseteq Q$.
	Set $T = \{  q \in Q \mid \exists u \in \Sigma^* \colon \delta(q, u) \in S \}$.
	Then $T$ is a weakly connected component of the automaton graph of~$\mathcal A$.
\end{corollary}
\begin{proof}
	Since the states in~$T$ are precisely those having a path leading into~$S$,
	and~$S$ is strongly connected, $T$ is weakly connected
	and thus contained in some weakly connected component.
	By the definition of~$T$, there are no transitions leading into~$T$
	since $\delta(q, a) \in T$ implies $q \in T$.
	Since~$T$ is an upper set by \cref{lem:comm:inv-images-letter-upper-set}
	there are also no transitions leading out of~$T$ since $\delta(q, a)$
	with $q \in T$ implies $\delta(q, a) \in T$.
	Thus~$T$ equals the weakly connected component
	containing it, since otherwise there would either exist a transition
	going out of~$T$ or into~$T$.
\end{proof}

\begin{corollary}  \label{cor:comm:inv-images-letter-upper-set}
	A commutative automaton $\semiautomat{A}{Q}{\Sigma}{\delta}$ whose automaton graph is weakly connected
	has a unique sink component.
\end{corollary}
\begin{proof}
	Let $S \subseteq Q$ be a sink component
	and consider $T = \{  q \in Q \mid \exists u \in \Sigma^* \colon \delta(q, u) \in S \}$.
	Since~$\mathcal A$ is weakly connected, by \cref{cor:comm:inv-images-letter-upper-set-T-weakly},
	we find $Q = T$.
	From \cref{lem:comm:inv-images-letter-upper-set} 
	we get that~$S$ is the only sink component in $\mathcal A$.
\end{proof}

\begin{proposition} \label{syncmax:classes:comm:charsync}	
	Let $\semiautomat{A}{Q}{\Sigma}{\delta}$ be a commutative automaton.
	Then the following are equivalent.
	\begin{enumerate}
		\item \label{syncmax:classes:comm:charsync:1} $\mathcal A$ has a trap state and its automaton graph is weakly connected.
		
		\item \label{syncmax:classes:comm:charsync:2} $\mathcal A$ has a trap state reachable from every other state.
		
		\item \label{syncmax:classes:comm:charsync:3} $\mathcal A$ is synchronizing.
	\end{enumerate}	
\end{proposition}
\begin{proof}
	$\ref{syncmax:classes:comm:charsync:1}\Rightarrow \ref{syncmax:classes:comm:charsync:2}$: By \cref{cor:comm:inv-images-letter-upper-set}
	the trap state forms a unique sink component.
	Thus, it is reachable from every other state.
	
	$\ref{syncmax:classes:comm:charsync:2}\Rightarrow \ref{syncmax:classes:comm:charsync:3}$: By \cref{prelim:thm:sync_trap_state}, $\mathcal A$
	is synchronizing.
	
	$\ref{syncmax:classes:comm:charsync:3}\Rightarrow \ref{syncmax:classes:comm:charsync:1}$: By \cref{constr_sync:prop:solvable_sync_iff_reach_trap},
	a synchronizing state is a trap state. Clearly, a synchronizing
	automaton has an automaton graph that is weakly connected.
\end{proof}

\Cref{syncmax:classes:comm:charsync} was stated in~\cite[Lem.~5.1]{DBLP:conf/fun/FernauH022}.

\subsection{Extremal Synchronizing Automata}
\label{prelim:classes:comm:subsec:extremal}

Since commutative automata are $\rho$-automata, and every synchronizing commutative automaton has a trap state, \cref{constr_input:prop:rho_extremal_sync_then_pto} implies that every synchronizing commutative $n$-state automaton with reset length $n-1$ is partially ordered. We now determine the structure of those commutative automata that are extremal w.r.t.~the reset length. Such automata were first classified in~\cite[Thm.~23]{FernauHoffmann19}. \Cref{fig:ex_extremal} shows a commutative automaton with the largest possible reset length.

\begin{figure}[ht]
	\centering
	\begin{tikzpicture}[>=latex',shorten >=1pt,node distance=4.5em, auto]
		
		\tikzset{
			state/.append style={
				inner sep=2pt, 
				minimum size=3pt,
			}
		}
		
		\node[state] (1) {$1$};
		\node[state] (2) [above right of=1] {$2$};
		\node[state] (3) [above right of=2] {$3$};
		\node[state] (4) [above right of=3] {$4$};
		\node[state] (5) [above left of=1] {$5$};
		\node[state] (6) [above left of=5] {$6$};
		
		\path[->] (1) edge [loop below] node {$a,b, c$} (1);
		\path[->] (2) edge  node {$b,c$} (1);
		\path[->] (3) edge  node {$b,c$} (2);
		\path[->] (4) edge  node {$b,c$} (3);
		\path[->] (2) edge [loop above] node {$a$} (2);
		\path[->] (3) edge [loop above] node {$a$} (3);
		\path[->] (4) edge [loop above] node {$a$} (4);
		
		\path[->] (5) edge [swap] node {$a$} (1);
		\path[->] (6) edge [swap] node {$a$} (5);
		\path[->] (5) edge [loop above] node {$b,c$} (5);
		\path[->] (6) edge [loop above] node {$b,c$} (6);
	\end{tikzpicture}	
	\caption[Extremal synchronizing commutative automaton]{Commutative six-state automaton with shortest synchronizing word $aabbb$.}
	\label{fig:ex_extremal}
\end{figure}

\begin{theorem}
	Let $\semiautomat{A}{Q}{\Sigma}{\delta}$ be a synchronizing commutative automaton
	with trap state~$t$.
	Then the reset length of~$\mathcal A$ equals $|Q| - 1$ if and only if 
	we have a partition $\Gamma_1, \ldots, \Gamma_m$ of $\Sigma$
	and sets $T_{\Gamma_i} \subseteq Q$ such that
	\begin{enumerate}
	\item $T_{\Gamma_i} \cap T_{\Gamma_j} = \{t\}$ for $i \ne j$ and $Q = T_1 \cup \cdots \cup T_m$,
	
	\item $T_{\Gamma_i} = \{ q_{i,0}, \ldots, q_{i, m_i}\}$ for some $m_i \ge 0$, % m_I = |T_{Gamma}|, mit zustand s_i?
	
	\item $\delta(q_{i,j}, x) = q_{i, \min\{j+1,m_i\}}$	for $x \in \Gamma_i$, and
	
	\item $\delta(q_{i,j}, x) = q_{i,j}$ for $x \in \Sigma \setminus \Gamma_i$.	
\end{enumerate}
\end{theorem}
\begin{proof}
That automata of this form have reset length $|Q| - 1$ is clear.	
	
Conversely, we do the proof only for a binary alphabet $\Sigma = \{a_1,a_2\}$, since for larger alphabets
the arguments are essentially the same, but involve more indices and more complicated notation.
% oder zwei die von q verzweigen q.a und q.b betrachten und erzeugtes upper set von {q.a,q.b} betrachten
% und die Fälle unterscheiden für beliebiges Alphabet

By \cref{constr_input:prop:rho_extremal_sync_then_pto}, $\mathcal A$ is partially ordered.
By \cref{prelim:general:order:cor:seq-upper-sets} there exists a sequence of upper sets
\[
S_1 \subsetneq S_2 \subsetneq \cdots \subsetneq S_m
\]
such that $S_1 = \{t\}$ and $S_i \setminus S_{i-1}$ for $i \in \{2,\ldots,m\}$ is a singleton set since~$\mathcal A$ is partially ordered. We show that the claim 
holds for each $\mathcal A_{|\mathcal S_i}$ by induction.

For $i = 1$, the claim clearly holds for the single-state automaton $\mathcal A_{|S_1}$.

Next, let $i \ge 2$. We have $S_i \setminus S_{i-1} = \{q\}$ for some $q \in Q$.
Without loss of generality, let $\delta(q, a_1) \ne q$.
By \cref{lem:sub-quotient-extremal-linear}, $\mathcal A_{|S_{i-1}}$
and $\mathcal A_{|S_i}$ have reset lengths $|S_{i-1}| - 1$ and $|S_i|-1$, respectively.
Thus, by induction hypothesis,
we either have $S_{i-1} = \{ q_0, \ldots, q_{l} \}$
with $q_l = t$ and $\delta(q_j, x) = q_{\min\{j+1,l\}}$ for $j \in \{0,1,\ldots,l\}$
for $x \in \Sigma$,
or $S_i = \{ q_{1, 0}, \ldots, q_{1, l_1}, q_{2, 0}, \ldots, q_{2, l_2}\}$
such that
\begin{enumerate}
\item $q_{l_1} = q_{l_2} = t$,

\item $\delta(q_{j,r}, a_j) = q_{j, \min\{r+1, l_j\}}$ for $j \in \{1,2\}$ and $r \in \{0,1,\ldots,l_j\}$,

\item $\delta(q_{j,r}, x) = q_{j,r}$ for $x \in \Sigma \setminus \{a_j\}$, $j \in \{1,2\}$ and $r \in \{0,1,\ldots,l_j\}$.
\end{enumerate}

First, suppose the case $S_{i-1} = \{ q_0, \ldots, q_{l} \}$ applies.
If $\delta(q, a_1) \in \{q_1, \ldots, q_l\}$, then $a_1^{l}$ synchronizes $\mathcal A_{|S_i}$,
which is not possible since $l = |S_{i-1}| - 1 < |S_i| - 1$.
Thus, $\delta(q, a_1) = q_0$. We also have $\delta(q, a_2) \notin \{q_1, \ldots, q_l\}$
since otherwise $a_2^l$ synchronizes~$\mathcal A_{|S_i}$.
Next, $\delta(q, a_2) = q$ is excluded, since then $a_2 a_1^{l}$ ends at the state $q_{l-1}$
and $a_1^la_2$ ends at $q_l$, which contradicts commutativity of $\mathcal A$.
Thus also $\delta(q, a_2) = q_0$. Then $\mathcal A_{|S_i}$ has the form as claimed.

Now, suppose the case $S_i = \{ q_{1, 0}, \ldots, q_{1, l_1}, q_{2, 0}, \ldots, q_{2, l_2}\}$
applies. Assume $\delta(q, a_1) \in \{q_{2,0},\ldots, q_{2,l_2}\}$
or $\delta(q, a_1) \in \{q_{1,1},\ldots,q_{1,l_1}\}$.
Then the word $a_1^{l_1} a_2^{l_2}$ also synchronizes $\mathcal A_{|S_i}$.
Thus, we must have $\delta(q, a_1) = q_{1,0}$.
Assume $\delta(q, a_2) \ne q$.
Then $\delta(q, a_2 a_1) \ne \delta(q, a_1a_2) = \delta(q, a_1)$,
contradicting the commutativity of~$\mathcal A$. Thus $\delta(q, a_2) = q$
and $\mathcal A_{|S_i}$ has the form as claimed.

Thus, we obtain the statement by induction.
\end{proof}